\documentclass[11pt,a4paper]{article}
\usepackage[utf8]{inputenc}
\usepackage{booktabs}
\usepackage{amsmath}
\usepackage{geometry}
\usepackage{hyperref}
\usepackage{xcolor}

\geometry{margin=2.5cm}

\title{Objective Function Analysis\\Basque Country Case Study}
\author{TransComp Model Documentation}
\date{January 2026}

\begin{document}

\maketitle

\section{Overview}

The optimization model's objective function value ($\sim$185.6 trillion EUR) is dominated by \textbf{budget constraint penalties}, not actual system costs. This document explains the breakdown and provides guidance for interpreting results.

\section{Objective Function Breakdown}

\subsection{Total Objective Value: $\sim$185.6 trillion EUR}

\begin{table}[h]
\centering
\begin{tabular}{lrr}
\toprule
\textbf{Component} & \textbf{Contribution} & \textbf{\% of Objective} \\
\midrule
Budget penalties (discounted) & $\sim$137 trillion EUR & $\sim$74\% \\
Other soft costs (VOT, LOS, etc.) & $\sim$48 trillion EUR & $\sim$26\% \\
\textbf{Real physical costs} & \textbf{$\sim$64 billion EUR} & \textbf{0.035\%} \\
\bottomrule
\end{tabular}
\caption{Objective function breakdown by component type}
\end{table}

\section{Real System Costs}

Discounted at 5\% annual rate from base year 2020.

\begin{table}[h]
\centering
\begin{tabular}{lrr}
\toprule
\textbf{Cost Component} & \textbf{Value} & \textbf{\% of Real Costs} \\
\midrule
Vehicle purchases & 55.4 billion EUR & 86\% \\
Maintenance & 5.8 billion EUR & 9\% \\
Energy/fuel & 3.1 billion EUR & 5\% \\
Infrastructure & 44 million EUR & $<$0.1\% \\
\midrule
\textbf{TOTAL} & \textbf{64.3 billion EUR} & 100\% \\
\bottomrule
\end{tabular}
\caption{Real system costs (discounted, 5\% rate)}
\end{table}

\section{Budget Penalty Mechanism}

\subsection{How It Works}

The model uses soft budget constraints with the formula:
\begin{equation}
\text{spending} \leq \text{budget} + \text{budget\_penalty\_plus} - \text{budget\_penalty\_minus}
\end{equation}

Where:
\begin{itemize}
    \item \texttt{budget\_penalty\_plus} = amount exceeding budget (penalized)
    \item \texttt{budget\_penalty\_minus} = amount under budget (also penalized)
    \item \textbf{Penalty coefficient = 10,000,000 EUR per EUR of violation}
\end{itemize}

\subsection{Why Such a Large Penalty?}

The 10M EUR penalty is a ``big-M'' numerical trick to:
\begin{enumerate}
    \item Make budget constraints nearly binding (households almost never exceed)
    \item Allow tiny violations for model feasibility
    \item Ensure the solver can find solutions even with tight budgets
\end{enumerate}

\subsection{Actual Budget Violations} 

\begin{table}[h]
\centering
\begin{tabular}{lrrr}
\toprule
\textbf{Income Group} & \textbf{Vehicles Purchased} & \textbf{Total Violation} & \textbf{Violation/Vehicle} \\
\midrule
Q1 (lowest) & 361,603 & 112M EUR & 310 EUR \\
Q2 & 463,089 & 129M EUR & 278 EUR \\
Q3 & 634,907 & 127M EUR & 200 EUR \\
Q4 & 992,555 & 13.5M EUR & 14 EUR \\
Q5 (highest) & 2,407,240 & 0 EUR & 0 EUR \\
\bottomrule
\end{tabular}
\caption{Budget violations by income quintile}
\end{table}

\textbf{Total actual budget excess: 381 million EUR} (0.27\% of vehicle spending)

In objective function: $381\text{M EUR} \times 10\text{M penalty} = 3{,}800$ trillion EUR (before discounting)

\section{Implications for Sensitivity Analysis}

\subsection{MIP Gap Interpretation}

With objective $\sim$185.6 trillion EUR and MIP gap of 0.01\%:
\begin{itemize}
    \item Threshold for significance: $\sim$18.5 billion EUR difference
    \item This is $\sim$300$\times$ the total infrastructure cost
    \item Small \% changes in objective are dominated by penalty changes, not real cost changes
\end{itemize}

\subsection{VOT Sensitivity Results}

\begin{table}[h]
\centering
\begin{tabular}{lr}
\toprule
\textbf{Metric} & \textbf{Value} \\
\midrule
Objective spread & $\sim$28 billion EUR (0.015\%) \\
Significance threshold & $\sim$37 billion EUR (0.02\%) \\
Result & Borderline significant \\
\bottomrule
\end{tabular}
\caption{VOT sensitivity analysis summary}
\end{table}

\subsection{Alpha (Circuity) Sensitivity Results}

\begin{table}[h]
\centering
\begin{tabular}{lr}
\toprule
\textbf{Metric} & \textbf{Value} \\
\midrule
Objective spread & $\sim$8 billion EUR (0.004\%) \\
Significance threshold & $\sim$18 billion EUR (0.01\%) \\
Result & NOT statistically significant \\
\bottomrule
\end{tabular}
\caption{Alpha sensitivity analysis summary}
\end{table}

\section{Recommendations}

\subsection{For Reporting Results}

\begin{enumerate}
    \item \textbf{Report real system costs separately} from objective value
    \item \textbf{Focus on distributional outcomes} (who buys EVs, who uses which chargers) rather than total cost differences
    \item \textbf{Frame sensitivity results} as showing model stability, not cost sensitivity
\end{enumerate}

\subsection{Suggested Framing}

\begin{quote}
``The sensitivity analysis reveals that total system costs are robust to variations in [parameter] ($<$0.02\% variation in objective), indicating model stability. However, distributional outcomes---particularly charging infrastructure access for lower-income groups---show meaningful variation across scenarios.''
\end{quote}

\subsection{For Future Model Runs}

Consider:
\begin{enumerate}
    \item Reducing penalty coefficient if budget constraint violations are acceptable
    \item Reporting a ``real cost'' metric alongside the objective
    \item Using tighter MIP gaps only when comparing scenarios with expected large differences
\end{enumerate}

\section{Technical Details}

\begin{table}[h]
\centering
\begin{tabular}{ll}
\toprule
\textbf{Parameter} & \textbf{Value} \\
\midrule
Discount rate & 5\% annual \\
Time horizon & 2020--2060 (40 years) \\
Investment periods & Every 5 years \\
\texttt{budget\_penalty\_plus} & 10,000,000 EUR \\
\texttt{budget\_penalty\_minus} & 10,000,000 EUR \\
\texttt{budget\_penalty\_yearly\_plus} & 10,000,000 EUR \\
\texttt{budget\_penalty\_yearly\_minus} & 10,000,000 EUR \\
\bottomrule
\end{tabular}
\caption{Key model parameters}
\end{table}

\vspace{1cm}
\noindent\rule{\textwidth}{0.4pt}

\noindent\textit{Generated: January 2026}\\
\noindent\textit{Case Study: cs\_concentrated\_revised\_2026-01-26\_00-48-33}

\end{document}
